\section{A logarithmic lower bound for the Poincaré constant}\label{sec:poincare}

The price comparison in Section~\ref{sec:comparison} requires a
single placement of the row sites that works for every balanced cut.
To construct it, we shall look for placements that make two random
changes of state have very different mean squared lengths. The
simplest comparison is on a regular graph. One step follows an edge,
and the other goes to an independent uniform vertex. If a class of
Hadamard spaces has no ports, the constant that bounds the second
energy by the first is at least of order \(\log n\) on \emph{every}
\(n\)-vertex graph of bounded degree. This bound will allow us to
distinguish two walks by the information they generate.

For a finite Markov kernel \(P\) with stationary law \(\pi\) and a
placement \(x=(x_i)\) of its states, put
\[
 \mathcal E_P(x)=\sum_{i,j}\pi_iP_{ij}d(x_i,x_j)^2,\qquad
 \mathcal V_\pi(x)=\sum_{i,j}\pi_i\pi_j d(x_i,x_j)^2 .
\]
For a pair law \(D\), the weights in \(\mathcal E_D(x)\) are its joint
probabilities. Write \(\mathcal V\) when \(\pi\) is uniform. States
of zero stationary mass may be discarded. A \(d\)-regular multigraph
has a symmetric matrix of nonnegative integral arc counts with every
row sum \(d\). A holding arc counts once. Dividing the counts by
\(d\) gives the walk kernel with uniform stationary law.

Suppose one large graph had too small a Poincar\'e constant. Placing
a copy of it at each vertex of a preceding graph would multiply the
order at every iteration. To keep the Poincar\'e constants bounded,
we average the walk on the preceding graph before inserting the
copies. The estimate for this averaging is a form of nonlinear
spectral calculus \cite{MendelNaor,MendelNaorBarycentric}. We give
the barycentric argument for Hadamard spaces
\cite{BH,Sturm}.

\begin{lemma}[Smoothing a walk]\label{lem:power}
Let \(P\) be a finite Markov kernel with stationary law \(\pi\), and let
\(t\ge2\) be even. Put
\[
 A=\tfrac12(I+P),\qquad S=\tfrac1t\sum_{h=1}^t A^h .
\]
For each placement \(x\) in a Hadamard space \(X\) there is a placement
\(y\) in \(X\) such that
\begin{equation}\label{eq:hadamard-smoothing}
 \sum_i\pi_i d(x_i,y_i)^2+t\mathcal E_P(y)
 \le 16\mathcal E_S(x).
\end{equation}
If \(D\) has both marginals \(\pi\) and
\(\mathcal E_D\le\beta\mathcal E_P\) for all placements in \(X\), then
\(\mathcal E_D\le C(1+\beta/t)\mathcal E_S\) for all placements in \(X\),
with an absolute constant \(C\).
\end{lemma}

\begin{theorem}[Logarithmic Poincaré obstruction]\label{thm:loggraph}
Let \(P\) be the walk on a \(d\)-regular multigraph with \(n\ge2\)
vertices, \(d\ge1\), and uniform stationary law \(\pi\). Suppose
that \(\mathcal V_\pi(x)\le\beta\mathcal E_P(x)\) for every
placement \(x\) in every member of a class of Hadamard spaces
without ports. Then
\[
                 \beta\ \ge\ c\,\frac{\log n}{\log(d+1)}
\]
for an absolute constant \(c>0\).
\end{theorem}

\begin{proof}[Proof of Lemma~\ref{lem:power}]
The barycenter \(b\) of a finite probability measure \(\mu\) in \(X\)
satisfies the variance inequality
\(\int d(z,w)^2\,d\mu(w)\ge d(z,b)^2+\int d(b,w)^2\,d\mu(w)\).
Indeed the midpoint inequality makes any minimizing sequence Cauchy,
and squared-distance convexity along the segment from the minimizer
to \(z\) gives the displayed inequality.
Starting with \(x^0=x\), take \(x_i^{g+1}\) to be the barycenter of the
\(x_j^g\) with weights \(A_{ij}\). Put
\[
 q_g=\sum_{i,j}\pi_iA_{ij}d(x_i^{g+1},x_j^g)^2,\qquad
 u_g=\sum_i\pi_i d(x_i,x_i^g)^2 .
\]
Run the stationary \(A\)-walk for \(h\) steps. Conditional on the first
\(h-1\) steps, apply the variance inequality to the last point with
the starting point as \(z\). Replace the last point by its conditional
barycenter and repeat. At the \(g\)-th replacement, starting at
\(g=0\), stationarity makes the mean variance term equal to \(q_g\).
Consequently
\[
 \mathcal E_{A^h}(x)\ge u_h+\sum_{g=0}^{h-1}q_g .
\]
On averaging over \(1\le h\le t\), each \(q_g\) with \(g<t/2\)
occurs in at least \(t/2\) summands. Since \(u_0=0\),
\[
 \sum_{0\le g<t/2}\bigl(u_g+(t/2)q_g\bigr)\le t\mathcal E_S(x).
\]
Choose a summand of size at most \(2\mathcal E_S(x)\).
Comparison along \(x_i^g,x_i^{g+1},x_j^g\) gives
\(\mathcal E_P(x^g)\le4q_g\), because \(A\) averages a holding step
and a \(P\)-step. Taking \(y=x^g\) proves \eqref{eq:hadamard-smoothing}.
Finally
\(\mathcal E_D(x)\le3(2u_g+\mathcal E_D(y))\) by the two marginal
identities. Apply the assumed comparison at \(y\).
\end{proof}

\begin{proof}[Proof of Theorem~\ref{thm:loggraph}]
A class without ports contains a space with two distinct
points. Choose one of the points independently and uniformly at
each vertex, and average the asserted inequality. The expected pair
energy is \((1-1/n)\) times half their squared distance, and the edge
energy is at most half their squared distance. Thus \(\beta\ge1/2\).
Write \(b=\max(1,\beta)\) and call the given multigraph \(H\).

Fix an even \(t\ge2\) with \((2d+1)^t\le n/2\). Given a previous
graph of degree \(2d\) with walk \(K\) and Poincaré constant at most
\(\gamma\), form the averaged kernel \(S\) of Lemma~\ref{lem:power}
from \(K\). The new pair comparison for \(S\) has constant at most
\(C(1+\gamma/t)\) throughout the class.
We need a symmetric approximation to \(S\) whose probabilities have
denominator \(n\). In each row set \(m_{uv}=\lceil n S_{uv}/2\rceil\).
A row of \(S\) has support at most \((2d+1)^t\), since holds pad shorter
walks to length \(t\). Hence its integer row sum is at most
\(n/2+(2d+1)^t\le n\). Pad each diagonal entry to reach \(n\).
The counts remain symmetric and \(R_{uv}=m_{uv}/n\ge S_{uv}/2\).
It follows that the Poincaré constant of \(R\) is at most
\(\rho=2C(1+\gamma/t)\).

Replace each old vertex \(u\) by a copy of \(H\). Partition its \(n\)
vertices into sets of sizes \(m_{uv}\), and match the set indexed by
\(v\ne u\) bijectively to the set indexed by \(u\) in the copy at \(v\).
Fix the vertices in the set indexed by \(u\) itself. This gives an
involution whose projected step from a uniform vertex of the copy
has distribution \(R_{u\bullet}\).
The new graph has the \(d\) internal arcs of \(H\) at each vertex and
\(d\) copies of its involution arc. Its mean edge energy is
\((E_{\rm in}+E_{\rm match})/2\), with the two terms the corresponding
mean energies.

For a placement \(x\) in any one of the test spaces, let \(z_u\)
be the barycenter of the points in the copy at \(u\). Let \(v_0\) be
the mean squared distance to one's own barycenter. Minimization against
a uniformly chosen point in each copy and the comparison on \(H\)
give \(v_0\le b E_{\rm in}\). Apply the triangle inequality first from
one barycenter to the next via a matched pair, then between independent
vertices via their two barycenters. Uniformity of both endpoints of an
involution step gives
\[
 \mathcal E_R(z)\le3(2v_0+E_{\rm match}),\qquad
 \mathcal V(x)\le3(2v_0+\mathcal V(z)).
\]
Here each pair energy uses the uniform law on the indicated graph.
Since \(\mathcal V(z)\le\rho\mathcal E_R(z)\), the new graph has a
Poincaré bound \(B b(1+\gamma/t)\) for an absolute \(B\ge1\).

If \(Bb/t\le1/2\), iteration starting from one vertex with \(2d\)
loops gives graphs on \(n^k\) vertices of fixed degree \(2d\) with
bounds at most \(2Bb\). These graphs would be ports for the class.
Hence every admissible even \(t\) satisfies \(t<2Bb\). If
\(\log(n/2)/\log(2d+1)\ge4\), the largest admissible even \(t\)
is at least half this ratio. In the remaining range
\(\log n/\log(d+1)\) is bounded absolutely and \(\beta\ge1/2\)
suffices. Since \(b\le2\beta\), this proves the assertion.
\end{proof}

We shall apply the bound to walks with unequal transition
probabilities. Only the number of possible destinations matters.

\begin{corollary}[Symmetric kernels with sparse rows]\label{cor:support}
Suppose \(P\) is a symmetric stochastic matrix on \(n\ge2\) states
with at most \(d\) positive entries per row, where \(d\ge1\) is
an integer, and let \(\pi\) be uniform. If
\(\mathcal V_\pi(x)\le\beta\mathcal E_P(x)\) for every placement \(x\)
in every member of a class of Hadamard spaces without ports, then
\(\log n\le C\beta\log(d+1)\) for an absolute constant \(C\).
\end{corollary}
\begin{proof}
Take \(\lceil dP_{ij}\rceil\) arcs in each position and add holding
arcs to make every row sum \(2d\). The walk on the resulting symmetric
regular multigraph dominates \(P/2\) entrywise, so its Poincaré
constant is at most \(2\beta\). Apply Theorem~\ref{thm:loggraph}.
\end{proof}
\section{Observable entropy}\label{sec:entropy}

A graph of small degree gives a walker few choices at each step.
A walk with many possible destinations can also transmit little
information if a few choices carry most of the probability, or if
many destinations share the same observable label. We compare
two walks on the same state space. One, the \emph{demand} walk,
specifies the pairs of states we want far apart in a placement.
The other, the \emph{physical} walk, specifies the moves we want
short. A uniform comparison of their squared lengths bounds the
amount of information one demand step can generate about the
destination label, given the full starting state. The bound measures
the physical walk by the number of bits per observed label along a
long stationary trajectory.

Given a labeling $L$ of a finite set $V$ and a stationary $P$-walk
$(X_t)_{t\ge0}$ on $V$, put
\[
 H_j=H\big(L(X_0),\ldots,L(X_{j-1})\big),\qquad
 a=\lim_{j\to\infty}\frac{H_j}{j}=\inf_{j\ge1}\frac{H_j}{j}.
\]
Here \(H\) denotes Shannon entropy in bits, and all logarithms in this
section are in base two \cite{ct}. The equality follows from
stationarity and subadditivity of block entropy.

\begin{theorem}[Observable entropy]\label{thm:entropy}
Let $P,D$ be Markov kernels on a finite set with common stationary
law $\pi$. Suppose that $\beta\ge1$ and that
$\mathcal E_D(f)\le\beta\mathcal E_P(f)$ for every placement $f$ in
a class of Hadamard spaces without ports. For any labeling $L$, let $a$ be its entropy rate
for the stationary $P$-walk and draw $(X,Y)$ from the stationary
$D$-step. Then
\[
 H(L(Y)\mid X)\le C\beta(1+a)
\]
for an absolute constant $C$.
\end{theorem}

In particular, if the demand information exceeds \(C\beta(1+a)\)
for some \(\beta\ge1\), one of the target spaces admits a placement with
\(\mathcal E_D(f)>\beta\mathcal E_P(f)\). Section~\ref{sec:product}
uses the identity labeling on a walk of words modulo cyclic rotation.
The full state is then observable and the entropy rate is its
one-step conditional entropy. In Section~\ref{sec:lamps} the label
records only an array of lamps. The physical walk also keeps auxiliary
records, so the cost in Theorem~\ref{thm:entropy} is the entropy rate
of the observed arrays over time.

The additive constant accounts even for deterministic motion. If \(P\)
alternates between two states and \(D\) draws a uniform state
independently, the identity labeling gives \(H_j=1\) for every \(j\).
Here \(a=0\), yet \(H(Y\mid X)=1\) and
\(\mathcal E_D=\mathcal E_P/2\).

\begin{proof}
We may assume that the stationary $D$-step is exchangeable. Indeed,
replace its pair law by the average of that law and its transpose.
The squared energies are unchanged, and concavity of conditional
entropy in the joint law shows that the new $H(L(Y)\mid X)$ is at
least half the original value. The stationary $P$-walk is unchanged.

Write $h=H(L(Y)\mid X)$ for this symmetrized step and assume $h>0$.
We use \(r\) copies of the original state, called the reservoir, to
write entries in an array of \(m\) lamps. At each writing step the
value added to one lamp is a function of the reservoir labels.
From a set of starts of mass at least one half, a demand move in all
\(r\) cells yields an increment whose law dominates a fixed fraction
of the uniform law on \(\lfloor rh/4\rfloor\)-bit words. We shall use the assumed
comparison to charge a refresh of all the lamps to \(O(mr\beta)\)
physical steps. Corollary~\ref{cor:support} will then require the net
change over those steps to carry order \(mrh\) bits. A description
of the same change using only the observed reservoir labels and the
writing schedule gives the upper bound in terms of \(a\).

\smallskip\noindent
\emph{Extracting lamp values.}
Fix \(m\ge1\). Take \(r\) sufficiently large that \(r\ge2\) and
\(k=\lfloor rh/4\rfloor\ge1\).
Set $A=\mathbb F_2^k$ and $q=|A|$. Write $L^r$ for the coordinatewise
labeling. There is a function $b$ from $r$-tuples of labels to $A$
and a set $B\subseteq V^r$ of $\pi^r$-mass at least $1/2$ such that
for $U'\sim D^{\otimes r}(u,\cdot)$ one has
\begin{equation}\label{eq:lamp-hash}
 \Pr\{b(L^r(U'))=z\mid U=u\}\ge \frac1{4q}
 \qquad(u\in B,\ z\in A).
\end{equation}
To see this, draw $r$ independent stationary $D$-steps. The negative
logarithm of the conditional probability of their destination label
tuple, given the starting reservoir $U$, is the sum of $r$ independent
random variables of mean $h$. The law of large numbers implies that,
for a set of starts of mass $1-o(1)$, atoms of conditional probability
at most $2^{-3rh/4}$ carry at least half the conditional mass. Assign
each label tuple a uniformly random color in $A$, independently. At
any such start, the mass of these small atoms assigned to a given
color has mean at least $1/(2q)$ and variance at most $2^{-3rh/4}/q$.
Chebyshev's inequality and a union bound give failure probability at
most $16q^2 2^{-3rh/4}=o(1)$ for the lower bound $1/(4q)$ in all
colors. Averaging over the coloring gives~\eqref{eq:lamp-hash}.

\smallskip\noindent
\emph{Comparing two walks on lamps and reservoirs.}
The physical kernel $G$ has stationary law $\pi^r$ times the uniform
law on $A^m$. With probability $1-1/r$ it advances a uniformly chosen
reservoir cell by $P$. With the remaining probability it chooses a
uniform lamp $J$ and adds $b(L^r(U))$ to that lamp. The first operation
preserves $\pi^r$ by stationarity of \(P\), and the second permutes the
uniform array for every fixed reservoir.
The demand kernel $D_*$ keeps the reservoir fixed. It replaces all
lamps by independent uniform values if $U\in B$ and holds otherwise.
For every placement \(F\) of this enlarged state space in a member of
the class, we claim that
\begin{equation}\label{eq:lamp-compare}
 \mathcal E_{D_*}(F)\le Cmr\beta\,\mathcal E_G(F).
\end{equation}

We first compare a refresh of just one lamp with the three moves in
Figure~\ref{fig:lamp}. Write \(s\in A^m\) for the lamp array and
\(e_Jz\) for the array equal to \(z\) in position \(J\) and zero
elsewhere. At a good starting reservoir \(u\in B\), draw
\(u'\sim D^{\otimes r}(u,\cdot)\). Add \(z=b(L^r(u'))\) to lamp \(J\) and return the reservoir to \(u\).
The endpoint is \((u,s+e_Jz)\). By \eqref{eq:lamp-hash}, every
increment has probability at least \(1/(4q)\). Thus the endpoint
law dominates one quarter of the uniform replacement law.

\begin{figure}[tbp]
\centering
\begin{tikzpicture}[x=1cm,y=1cm,font=\small,
    state/.style={rectangle,rounded corners=2pt,draw=ink,
                  fill=ink!3,inner xsep=5pt,minimum width=2.55cm,
                  minimum height=7.5mm},
    move/.style={-{Stealth[length=2mm]},semithick,draw=ink}]
  \node[state] (ul) at (0,1.7) {\((u,s)\)};
  \node[state] (ur) at (7.6,1.7) {\((u,s+e_Jz)\)};
  \node[state] (ll) at (0,0) {\((u',s)\)};
  \node[state] (lr) at (7.6,0) {\((u',s+e_Jz)\)};
  \draw[move,dashed] (ul.east) -- node[above] {refresh lamp \(J\)} (ur.west);
  \draw[move] (ul.south) -- node[right] {draw \(u'\) by \(D^{\otimes r}\)} (ll.north);
  \draw[move,draw=sea] (ll.east) -- node[below] {add \(z=b(L^r(u'))\) to lamp \(J\)} (lr.west);
  \draw[move] (lr.north) -- node[left] {restore \(u\)} (ur.south);
\end{tikzpicture}
\caption{Refreshing lamp \(J\) while restoring the reservoir to its
initial state \(u\). The dashed step adds a uniform increment. On the
solid path the vertical moves are charged to the demand walk on the
reservoir cells, and the horizontal move is a single physical lamp
addition. For \(u\in B\), each endpoint of the dashed step occurs
along the solid path with at least one quarter of its probability.}
\label{fig:lamp}
\end{figure}

To pay for the outward reservoir move, we need the squared cost of
changing independent cells to be bounded by the sum of their
individual costs. If $(S_i,S'_i)$ are independent exchangeable
pairs and $S^{(i)}$ denotes $S$ with $S_i$ replaced by $S'_i$, then
\begin{equation}\label{eq:product-exchange}
 \mathbb E\,d(F(S),F(S'))^2
 \le\sum_i\mathbb E\,d(F(S),F(S^{(i)}))^2.
\end{equation}
For two blocks of coordinates form the quadrilateral obtained by
choosing the original or changed value in each block. The CAT$(0)$
quadrilateral inequality bounds the sum of squared diagonals by the
sum of squared sides. Exchanging one block preserves the joint law,
so the diagonals have equal expected squares, and each opposite pair
of sides does too. Dividing by two and iterating over the blocks
proves~\eqref{eq:product-exchange}.

Make independent $D$-moves in all $r$ reservoir cells, keeping the
lamps fixed. Under stationarity the squared cost of this change is at
most $2r\beta\,\mathcal E_G(F)$ by~\eqref{eq:product-exchange} and
the assumed comparison on each cell, since $1-1/r\ge1/2$. The squared
cost of a single random lamp addition is at most $r\mathcal E_G(F)$.
It remains to charge the return move. Integrate the path bound over
\(u\in B\), a uniform lamp array \(s\), and a uniform address \(J\).
The squared endpoint distance is at most three times the sum of the
squared lengths of the three legs. We bound this sum by integrating
over all starts. The first and last legs then have equal average cost,
since $s+e_Jz$ is uniform even conditional on $u,u',J$. The middle leg
has the stationary lamp-addition cost because \(u'\) has law \(\pi^r\).
Thus replacing one random lamp at good starts, with holding elsewhere,
costs at most $Cr\beta\,\mathcal E_G(F)$. Applying
\eqref{eq:product-exchange} to the $m$ independent pairs of old and
new lamp values at each good start gives~\eqref{eq:lamp-compare}.

\smallskip\noindent
\emph{Entropy forced by spreading the lamps.}
Take $T=2\lceil mr\beta\rceil$ and apply Lemma~\ref{lem:power}
to~\eqref{eq:lamp-compare}, using the Ces\`aro average of the lazy
$G$-walk up to time $T$. For placements depending only on the lamps,
the demand energy is $\pi^r(B)\mathcal V(f)$, where $\mathcal V(f)$
is the mean squared distance between images of two independent
uniform elements of $A^m$. The physical displacement of the lamps is
a random translation by some $Z\in A^m$, independent of their uniform
starting value $s$. All reservoir moves and choices of lamps are
independent of $s$. We obtain, with an absolute $K\ge1$,
\begin{equation}\label{eq:lamp-translation}
 \mathcal V(f)\le K\,\mathbb E\, d(f(s),f(s+Z))^2
 \qquad\text{for all placements } f \text{ in the class} .
\end{equation}

Each fixed translation has squared energy at most $4\mathcal V(f)$,
by the triangle inequality through an independent uniform point.
Put $\eta=1/(8K)$. Replace by zero every value of $Z$ whose probability
is less than $2^{-H(Z)/\eta}$. The replaced mass is at most $\eta$
by Markov's inequality for the information of the sampled atom (if
$H(Z)=0$ the support is a singleton). The discarded translations
contributed at most
\(4\eta\mathcal V(f)\) to the expectation in
\eqref{eq:lamp-translation}. Since \(4K\eta=1/2\), the resulting
translation kernel satisfies the same comparison with constant \(2K\).
It is symmetric because $A^m$ has exponent two, and each of its rows
has at most $1+2^{H(Z)/\eta}$ atoms. Corollary~\ref{cor:support}
therefore gives
\begin{equation}\label{eq:lamp-bits}
 mk=\log |A^m|\le C(1+H(Z)).
\end{equation}

\smallskip\noindent
\emph{Coding the physical displacement.}
In a trajectory of the lazy Ces\`aro walk, ignore lazy holds. There
are at most $T$ steps left. Let $S$ be the number of lamp additions,
so $\mathbb E S\le T/r$. Record $S$ in unary and the $S$ lamp
addresses in $\lceil\log m\rceil$ bits each. Within each of the $S+1$
intervals before, between and after additions, record in unary the
number of $P$-steps taken by each reservoir cell. The lengths of
these $r(S+1)$ unary words sum to at most $T+r(S+1)$ bits.
This prefix code for the schedule data has expected length at most
\[
 C\big(T+r+(T/r)\log(2m)\big).
\]
The order of moves in different cells inside an interval is
immaterial, since no lamp is read or written there.

Given the recorded data, the successive labels in a cell have the law
of a stationary $P$-block of the indicated length $j$, including its
initial label. This follows by constructing independent stationary
$P$-trajectories for the cells before sampling the schedule. Since
$H_j/j\to a$, there is $M_0<\infty$, depending only on $P,\pi,L$,
such that $H_j\le (1+a)j+M_0$ for every $j$. The $r$ lengths sum to
at most $T+r$. The schedule and these label blocks determine the
reservoir labels at every lamp addition and hence determine $Z$.
The chain rule for entropy now gives
\begin{equation}\label{eq:lamp-code}
 H(Z)\le C\big((T+r)(1+a)+rM_0+(T/r)\log(2m)\big).
\end{equation}
Combine~\eqref{eq:lamp-bits} and~\eqref{eq:lamp-code} and divide by
$mr$. Let $r\to\infty$ with $m$ fixed. The lamp-address contribution
tends to zero, and $k/r\to h/4$. Next let $m\to\infty$ to remove
$M_0/m$. As $T/(mr)\to 2\beta$ in the first limit, this proves the
theorem.
\end{proof}

\subsection{Conditional entropy of weighted edge laws}\label{sec:cursor-entropy}

For a graph with weighted edges, the demanded pairs and the
physical edges can have different vertex marginals. Symmetrize each
pair law and add suitable holding steps. The resulting walks have a
common stationary law and both squared energies are halved. Applying Theorem~\ref{thm:entropy} with the
identity labeling then bounds the one-step demand entropy.

\begin{corollary}[Pair laws on a graph]\label{cor:graph-entropy}
Let $\mu$ and $\alpha$ be laws on ordered pairs from a finite set, and
suppose that $\mathcal E_\mu(f)\le\beta\mathcal E_\alpha(f)$
for every placement \(f\) in every member of a class of Hadamard
spaces without ports, with \(\beta\ge0\). Put
$\alpha_s=(\alpha+\alpha^\top)/2$. Then
\[
 H_\mu(Y\mid X)\le C(1+\beta)\big(1+H_{\alpha_s}(Y\mid X)\big).
\]
In particular, if the undirected support of $\alpha$ has at most
$d$ neighbors at each vertex, $d\ge2$, then
$H_\mu(Y\mid X)\le C'(1+\beta)\log d$.
\end{corollary}
\begin{proof}
Put $\mu_s=(\mu+\mu^\top)/2$, and write $\rho,\sigma$ for the
marginals of $\mu_s,\alpha_s$, respectively. Write $\Delta_\rho$ for the
law of $(X,X)$ with $X\sim\rho$. The stationary pair laws
$(\mu_s+\Delta_\sigma)/2$ and $(\alpha_s+\Delta_\rho)/2$ define
kernels $D,P$ with common stationary law $(\rho+\sigma)/2$ on its
support. Their energies are half the original energies. Concavity of
conditional entropy gives $H_D(Y\mid X)\ge H_\mu(Y\mid X)/4$.
For the physical step, record a bit indicating which component of the
mixture was used. Conditional on the diagonal component the destination
is known, so $H_P(Y\mid X)\le 1+H_{\alpha_s}(Y\mid X)/2$.
This conditional entropy is the entropy rate for the identity labeling
of the stationary $P$-chain. The first assertion follows by applying
Theorem~\ref{thm:entropy} with $\max(1,\beta)$. In the bounded-degree
case at most $d$ neighbors and a possible holding move give
$H_{\alpha_s}(Y\mid X)\le\log(d+1)\le2\log d$.
\end{proof}
